\documentclass[12pt]{article}
\usepackage[margin=1in]{geometry}
\usepackage{setspace}
\usepackage{enumitem}
\onehalfspacing

\begin{document}
\thispagestyle{empty}

\begin{center}
{\large \textbf{Key Public Health Message}}
\end{center}

\vspace{0.5cm}

\textbf{Title:} Estimating Infectious Disease Importation Risk during the 2026 FIFA World Cup in the United States, June--July 2026

\vspace{0.5cm}

\begin{itemize}[leftmargin=*, itemsep=6pt]

  \item The 2026 FIFA World Cup will concentrate visitors from 48 nations across 11 US host cities between June 11 and July 19, creating conditions for accelerated importation of multiple infectious diseases simultaneously.

  \item Using a Poisson importation framework with Monte Carlo uncertainty propagation, we estimate that dengue fever poses the highest importation risk at most venue cities, with more than 10 expected importations from Brazil alone; malaria exceeds dengue probability at Atlanta, driven by direct travel from West and Central Africa.

  \item Background tourism — travel that would occur regardless of the tournament — accounts for the dominant share of total importation risk; the World Cup fan increment contributes approximately 8.3\% of projected arrivals for qualified nations, underscoring that surveillance should not focus exclusively on fan travel.

  \item The framework is built entirely from publicly available data (WHO surveillance, BTS T-100, CBP I-94), is fully reproducible, and is immediately extensible to future mass gathering events and additional pathogens.

\end{itemize}

\end{document}
